\section{Verwandte Arbeiten}
\label{sec:rw}
Dieses Kapitel stellt die wissenschaftlichen Grundlagen und den Forschungsstand in den für diese Arbeit relevanten Bereichen dar. Zunächst wird das Prinzip der körperbasierten Interaktion eingeführt. Darauf aufbauend werden immersive projektionsbasierte Interaktionsumgebungen und deren spezifische Anforderungen an die Interaktionsgestaltung, sowie controllerfreie Mid-Air- und Whole-Body-Interaktion betrachtet. Anschließend werden Navigation in immersiven Umgebungen, Gestendesign, Auswahlmechanismen und menübasierte Interaktion diskutiert. Abschließend wird der Kontext multisensorischer und atmosphärischer Interaktion hergestellt, um die konzeptionelle Einordnung dieser Arbeit und die zugrunde liegende Forschungslücke zu verdeutlichen.

Ziel ist es, bestehende Ansätze hinsichtlich ihrer Eignung für eine controllerfreie, körperbasierte Interaktion in immersiven, projektionsbasierten und multisensorischen Umgebungen zu analysieren und daraus eine Forschungslücke abzuleiten, die durch den MIR und die in dieser Arbeit entwickelten Interaktionskonzepte adressiert wird.



\subsection{Körperbasierte Interaktion (Embodied Interaction)}
\label{sec:rw_embodied}

Körperbasierte Interaktion (Embodied Interaction) nutzt den menschlichen Körper, seine Bewegungsfähigkeiten und sein Verständnis der physischen Welt als zentrale Grundlage der Interaktion mit digitalen Systemen \parencite{jacob2008realitybasedinteraction}. Dieses Paradigma baut auf \textcite{jacob2008realitybasedinteraction}s Prinzip der ``Reality-Based Interaction'' (RBI) auf, das die natürliche Interaktion des Menschen mit der realen Welt als Vorbild für die Gestaltung von Benutzungsschnittstellen nutzt.

Ziel der körperbasierten Interaktion ist es, die kognitive Übersetzung zwischen Intention und Systemreaktion zu reduzieren. \textcite{norman1986CognitiveEngineering} beschreibt diese Distanz als ``Gulf of Execution''. Werden Interaktionen an reale Handlungen angelehnt, kann der mentale Aufwand reduziert werden, der erforderlich ist, um eine gewünschte Wirkung im System zu erzielen \parencite{norman1986CognitiveEngineering}. Im Unterschied dazu erfordern controllerbasierte Eingaben häufig eine zusätzliche mentale Abbildung zwischen physischer Handlung und digitaler Reaktion.

Darüber hinaus kann die enge Kopplung von Körperbewegung und Systemfeedback in immersiven Umgebungen das Gefühl von Präsenz, also das Empfinden, tatsächlich ``vor Ort'' zu sein, verstärken \parencite{slatersanchesvivies2016enhancingourliveswithimmersivevirtualreality, slater1997frameworkForImmersiveVirtualEnvironments}. Studien zeigen, dass Nutzer:innen Gestensteuerung als immersiver, intuitiver und ansprechender empfinden als klassische controllerbasierte Eingaben \parencite{spittle2023reviewinteractiontechniquesforimmersiveenvironments, vona2025handsvscontrollerscomparing}.

Gleichzeitig bringt der Ansatz Herausforderungen mit sich. In der Literatur werden wiederholt Herausforderungen hinsichtlich Präzision, physischer Ermüdung und der Gestaltung verständlicher Interaktionsmetaphern beschrieben \parencite{spittle2023reviewinteractiontechniquesforimmersiveenvironments, kostic2024exploringmidairhandinteractionindatavisualization}.
\Textcite{jacob2008realitybasedinteraction} argumentieren entsprechend, dass realitätsnahe Interaktionen zugunsten von Effizienz, Ergonomie oder technische Umsetzbarkeit jedoch bewusst abstrahiert oder eingeschränkt werden können.

Aktuelle empirische Arbeiten bestätigen dieses Spannungsfeld. So zeigen \textcite{vona2025handsvscontrollerscomparing} in einer vergleichenden Studie zu handbasierter und controllerbasierter Interaktion in einer hochrealistischen VR-Shoppingumgebung, dass handbasierte Eingaben aufgrund ihrer Nähe zu realweltlichen Greifbewegungen als intuitiver und immersiver wahrgenommen werden, während controllerbasierte Eingaben Vorteile hinsichtlich Präzision und Zuverlässigkeit aufweisen \parencite{vona2025handsvscontrollerscomparing}. Sie beobachten zudem, dass das Fehlen haptischer Rückmeldung sowie eine reduzierte Greifgenauigkeit bei handbasierter Interaktion zu wiederholten Manipulationsversuchen und erhöhtem Frustrationspotenzial führen können. Insgesamt verdeutlichen diese Befunde, dass körperbasierte Interaktion zentrale Ziele der Reality-Based Interaction unterstützt, ihre Gestaltung jedoch insbesondere bei präzisionskritischen Aufgaben eine sorgfältige Abstimmung von Natürlichkeit, Feedback, Ergonomie und Robustheit erfordert.


\subsection{Immersive räumliche Interaktionsumgebungen und Projektionsräume}
\label{sec:rw_spatial}

Immersive Umgebungen zielen darauf ab, die Grenzen zwischen physischer und digitaler Welt aufzulösen, indem Nutzer:innen digitale Inhalte nicht nur betrachten, sondern den Raum selbst als Ort der Interaktion erleben \parencite{slater1997frameworkForImmersiveVirtualEnvironments}. Laut \textcite{slater1997frameworkForImmersiveVirtualEnvironments} ist ein zentraler Faktor dafür die enge Verknüpfung von Wahrnehmung und körperlicher Bewegung. Interaktion wird dabei nicht mehr als isolierte Eingabehandlung verstanden, sondern implizit von natürlichen Bewegungen der Nutzer:innen abgeleitet. % als Teil eines kontinuierlichen Zusammenspiels von Körper, Raum und Systemreaktion.

Obwohl \textcite{slater1997frameworkForImmersiveVirtualEnvironments} ihre Überlegungen im Kontext virtueller Umgebungen entwickeln, lassen sich die beschriebenen Prinzipien körperlich gebundener Wahrnehmung und räumlicher Interaktion auch auf andere immersive oder räumliche digitale Umgebungen übertragen.
Neben Head-Mounted-Display-basierten XR-Systemen zählen hierzu insbesondere projektionsbasierte Installationen, Medienräume oder interaktive Umgebungen, in denen der physische Raum selbst zur Interaktionsmöglichkeit wird. Charakteristisch ist die Abkehr von klassischen Desktop-Metaphern hin zu Interaktionsformen, die Körper, Raum und Kontext als integrale Bestandteile der Eingabe und des Feedbacks nutzen.

Zur Beschreibung räumlicher Interaktion ist zudem eine Einordnung nach grundlegenden Interaktionsaufgaben hilfreich. In räumlichen Benutzungsoberflächen lassen sich insbesondere die Aufgaben Selektion, Manipulation und Navigation unterscheiden \parencite[Kap.~5]{LaViola20043DUserInterfaces}. Selektion bezeichnet dabei die Auswahl von Objekten oder Funktionen, Manipulation deren Veränderung oder direkte Modifizierung und Navigation die Fortbewegung innerhalb einer virtuellen oder projizierten Umgebung. Diese Unterscheidung bietet eine nützliche Strukturierung, um unterschiedliche Interaktionsansätze in räumlichen Systemen vergleichend einzuordnen.\\

Das Konzept der ``Proxemic Interaction'' nach \textcite{ballendatmarquardtgreenberg2010proxemicinteraction} knüpft daran an und beschreibt, wie räumliche Beziehungen wie Distanz, Orientierung und Bewegung zwischen Nutzer:innen und Systemen gezielt als Interaktionsparameter genutzt werden können. %Interaktion entsteht dabei sowohl implizit, etwa durch Annäherung oder Blickrichtung, als auch explizit durch gezielte Handlungen im Raum.
Diese Idee findet sich auch in projektionsbasierten Umgebungen wieder: So reagieren Systeme wie ``RoomAlive'' oder ``LightSpace'' nicht nur auf explizite Eingaben, sondern auf Position, Annäherung und Bewegung im Raum.
Sie demonstrieren in diesem Zuge, wie die gesamte Raumarchitektur durch eine Kombination aus Projektoren und Tiefenkameras zu einer interaktiven Fläche werden kann \parencite{jones2014roomalive}. ``RoomAlive'' \parencite{jones2014roomalive} transformiert mittels Projektor- und Tiefenkamerasystemen einen gesamten Raum in eine augmentierte Umgebung, in der Nutzer:innen durch Ausweichen, Werfen oder Berühren körperlich mit digitalen Inhalten interagieren. Ähnlich erweitert ``LightSpace'' \parencite{wilsonbenko2010lightspace} die Interaktion über mehrere physische Oberflächen hinweg und ermöglicht Eingaben auf, über und zwischen physischen Strukturen.
Gleichzeitig zeigen diese Systeme typische Herausforderungen projektionsbasierter Interaktionsräume, darunter Kalibrierungsfehler, Trackingprobleme bei überlappenden Sensorsystemen, Systemlatenzen sowie die Schwierigkeit, Aufmerksamkeit innerhalb komplexer räumlicher Szenarien gezielt zu lenken \parencite{jones2014roomalive, wilsonbenko2010lightspace}.
Trotz dieser Herausforderungen, verdeutlichen beide Systeme den Übergang von gerätezentrierter Bedienung hin zu Ganzkörperinteraktion innerhalb einer augmentierten räumlichen Umgebung. \\

%{\color{red} Large displays wirklich wichtig oder lieber raus?}
%Neben solchen immersiven Installationen wurde räumliche Interaktion auch im Kontext großformatiger Displays untersucht. Large High-Resolution Display Environments (LHRDs) erweitern klassische Interaktionsparadigmen, indem sie physische Navigation, Multi-User-Koordination und die Nutzung verschiedener Eingabegeräte erfordern. Durch ihre Größe und Auflösung ermöglichen sie es, große Informationsmengen gleichzeitig darzustellen und fördern die Nutzung von Bewegung im Raum als Teil der Interaktion \parencite{belkacem2024LHRDsurvey}. Gleichzeitig entstehen neue Herausforderungen, etwa hinsichtlich der physischen Erreichbarkeit von Inhalten, möglicher Informationsüberlastung sowie der Koordination mehrerer Nutzer:innen. Zudem existieren keine einheitlichen Interaktionsstandards, sodass je nach Distanz und Kontext unterschiedliche Eingabemodalitäten wie Touch, Gesten, mobile Geräte oder Blicksteuerung zum Einsatz kommen \parencite{belkacem2024LHRDsurvey}.

% {\color{red} Oder das stattdessen erwähnen? Oder beides ganz raus?}
%Obwohl Large High-Resolution Display Environments (LHRDs) ähnliche Herausforderungen in Bezug auf physische Navigation und Multi-User-Interaktion aufwerfen \parencite{belkacem2024LHRDsurvey}, unterscheiden sie sich von projektionsbasierten multisensorischen Räumen wie dem MIR durch die fehlende Integration von nicht-visuellen Sinneskanälen und die geringere räumliche Immersion. Daher werden LHRD-Aspekte hier nur insofern relevant, als sie Hinweise auf die Bedeutung von Bewegung im Raum und räumlicher Koordination liefern.

Insbesondere in öffentlich zugänglichen oder gemeinsam genutzten Interaktionsumgebungen spielen darüber hinaus soziale Faktoren eine zentrale Rolle. \textcite{mueller2010requirementsAndDesignSpaceForInteractivePublicDisplays} zeigen, dass Interaktionen so gestaltet werden müssen, dass sie Aufmerksamkeit erzeugen, ohne Nutzer:innen unangenehm ausgesetzt darzustellen. Gleichzeitig sollten Systeme unterschiedliche soziale Rollen berücksichtigen, zwischen impliziten und expliziten Interaktionsformen wechseln können und Motivation durch Faktoren wie Neugier, Wahlmöglichkeiten oder Zusammenarbeit fördern.
Daraus ergibt sich für die Interaktionsgestaltung, dass Bewegungen sozial akzeptabel, körperlich verhältnismäßig und in gemeinsam genutzten Räumen ohne hohe Hemmschwelle ausführbar sein sollten.

Insgesamt verdeutlichen die betrachteten Arbeiten einen grundlegenden Wandel in der Gestaltung interaktiver Systeme: Interaktion verlagert sich zunehmend vom einzelnen Eingabegerät in den Raum selbst. Körperliche Bewegung, räumliche Position und Orientierung werden dabei zu zentralen Interaktionsressourcen. Gleichzeitig entstehen neue Anforderungen an die Interaktionsgestaltung, insbesondere hinsichtlich Ergonomie, sozialer Akzeptanz, räumlicher Einbettung und der Abstimmung von Interaktion und Wahrnehmung.

Diese Erkenntnisse sind insbesondere für projektionsbasierte multisensorische Umgebungen wie den MIR relevant. Sie machen deutlich, dass Interaktion in solchen Systemen nicht isoliert betrachtet werden kann, sondern stets im Zusammenspiel von Raum, Körper und Wahrnehmung gestaltet werden muss.


\subsection{Controllerfreie Mid-Air und Whole-Body Interaktion}
\label{sec:rw_midair}

Innerhalb solcher räumlichen Interaktionsumgebungen stellen Mid-Air- und Whole-Body-Interaktionen zentrale controllerfreie Eingabeformen dar. Anstelle physischer Eingabegeräte wie Maus, Tastatur oder Controller werden hierbei Handbewegungen, Gesten oder Körperpositionen direkt als Eingabeparameter genutzt \parencite{spittle2023reviewinteractiontechniquesforimmersiveenvironments, gavgiotaki2023gesturebasedinteractionforarsystemsreview}.
Aufbauend auf den Prinzipien körperbasierter Interaktion stellen Mid-Air- und Whole-Body-Interaktionen zentrale controllerfreie Eingabeformen in räumlichen Umgebungen dar.
Insbesondere in räumlichen und multisensorischen Umgebungen ermöglichen Mid-Air- und Whole-Body-Interaktionen eine enge Verbindung zwischen physischem Verhalten und digitalem Feedback und unterstützen damit zentrale Prinzipien körperbasierter Interaktion.

Gleichzeitig zeigt die Forschung, dass controllerfreie Interaktion mit spezifischen Herausforderungen verbunden ist, die bei der Gestaltung entsprechender Systeme berücksichtigt werden müssen.

Ein zentrales Problem stellt die physische Ermüdung dar. Das längere Halten der Arme im Raum führt zum sogenannten ``Gorilla-Arm''-Effekt, der die Nutzungsdauer und den Komfort erheblich einschränken kann \parencite{hincapi2014metrictoquantifyarmfatigue, spittle2023reviewinteractiontechniquesforimmersiveenvironments}. Studien zeigen, dass insbesondere statische Armhaltungen und große Bewegungsamplituden die Belastung erhöhen, während kompakte und körpernahe Interaktionen ergonomischer sind \parencite{hincapi2014metrictoquantifyarmfatigue}.

Darüber hinaus weisen Mid-Air-Interaktionen im Vergleich zu physischen Eingabegeräten häufig eine geringere Präzision auf. Das Fehlen taktiler Rückmeldung erschwert die präzise Ausführung und Bestätigung von Eingaben und kann zu Unsicherheiten, erhöhtem Fehleraufkommen sowie geringerem Vertrauen in das System führen \parencite{spittle2023reviewinteractiontechniquesforimmersiveenvironments, vona2025handsvscontrollerscomparing}. Dies ist insbesondere bei präzisionskritischen Aufgaben relevant.

Ein weiteres grundlegendes Problem ist die sogenannte Discoverability von Interaktionen. Da Gesten im Gegensatz zu grafischen Benutzungsoberflächen nicht sichtbar sind, müssen Nutzer:innen wissen oder erlernen, welche Interaktionen möglich sind und wie diese auszuführen sind \parencite{vogiatzidakis2018gestureeliticationstudiesreview}. In diesem Zusammenhang beschreiben \textcite{wobbrock2005discoverability} das Konzept der ``Guessability'', also die Fähigkeit, die Bedeutung einer Geste ohne vorherige Einweisung korrekt zu erschließen. Eine geringe Discoverability kann die Nutzbarkeit gestenbasierter Systeme erheblich einschränken. In der Literatur werden hierfür unter anderem visuelle Hinweise, Training oder menübasierte Unterstützungsmechanismen diskutiert \parencite{gavgiotaki2023gesturebasedinteractionforarsystemsreview}.

Eng damit verbunden ist das Problem unbeabsichtigter Eingaben, häufig als ``Midas Touch''-Problem bezeichnet. Alltägliche Bewegungen können vom System fälschlicherweise als Interaktion interpretiert werden, wenn keine klaren Mechanismen zur Trennung von Interaktion und Nicht-Interaktion vorhanden sind. Zur Adressierung dieses Problems werden in der Literatur unter anderem explizite Aktivierungsgesten zur Erreichung von quasi-modalen Interaktionszuständen vorgeschlagen, die die Möglichkeit zur Interaktion bewusst ein- und ausschalten, beispielsweise indem erst nach einer expliziten Aktivierungsgeste weitere Interaktionsgesten registriert werden.
\parencite{walterBaillyValkanova2014CuenesicsUsingMidAirGesturesToSelectItemsOnInteractivePublicDisplays}.

%{\color{red} Das lieber in Gestenabschnitt?}
%Zusätzlich zeigt sich beim Design von Gestensets, was in Elicitation Studien, in denen Nutzer:innen selbst Gestensets definieren sollen, dass Nutzer:innen dazu neigen, bekannte Interaktionsmuster aus Desktop- oder Touch-Interfaces auf den dreidimensionalen Raum zu übertragen. Dieses als ``Legacy Bias'' bezeichnete Verhalten kann zu ineffizienten oder unergonomischen Interaktionsformen führen \parencite{pasquale2024empiricalevaluationdefiningmidair}, hat jedoch auch das Potenzial der verbesserten Memorisierbarkeit. % Gleichzeitig besteht eine starke Präferenz für pantomimische Gesten, die reale Handlungen wie Greifen oder Zeigen nachahmen, da diese als besonders intuitiv wahrgenommen werden \parencite{kostic2024exploringmidairhandinteractionindatavisualization, gavgiotaki2023gesturebasedinteractionforarsystemsreview}. {\color{red} Das hier vielleciht lieber weg, weil das unten ja erstmal in Gesten Abschnitt richtig behandelt wird?}

Insgesamt zeigt die Literatur, dass Mid-Air- und Whole-Body-Interaktion ein hohes Potenzial für natürliche und immersive Interaktion bieten, gleichzeitig jedoch eine sorgfältige Gestaltung erfordern. Insbesondere müssen Aspekte wie Ergonomie, Präzision, Rückmeldung, Discoverability und die Vermeidung unbeabsichtigter Eingaben berücksichtigt werden. Diese Herausforderungen verdeutlichen, dass controllerfreie Interaktion nicht allein durch die Nutzung natürlicher Bewegungen entsteht, sondern durch eine gezielte Kombination aus geeigneten Interaktionsmechanismen, klarer Rückmeldung und strukturierten Interaktionskonzepten gestaltet werden muss.

% Um diese Grenzen zu überwinden, rücken multimodale Ansätze zunehmend in den Fokus: Evangelou et al. (2024)(QUELLE) zeigen, dass die Kombination von Mid-Air-Gesten mit haptischem Feedback die Wahrnehmung von Kontrolle und Vertrauen erhöht, während Vieira et al. (2024) ein Design-Space-Modell vorschlagen, das die konsistente Integration von Gesten, Sprache und Feedback als zentral für glaubwürdige Interaktionen identifiziert. Ein weiterer Trend ist die Entwicklung whole-body-interaktion, die den Fokus von den Händen auf den gesamten Körper ausdehnt: Position, Haltung und Bewegung als zentrale Interaktionsmittel (Matar et al., 2023).


\subsection{Navigation in immersiven Umgebungen}
\label{sec:rw_navigation}
Neben Selektion und Manipulation stellt auch Navigation eine grundlegende Interaktionsaufgabe in immersiven Umgebungen dar, da Nutzer:innen sich innerhalb virtueller oder projizierter Räume, die aber nur einen begrenzten physischen Raum enthalten, orientieren und fortbewegen können sollen. In der Forschung zu immersiven Interaktionssystemen werden unterschiedliche Lokomotionsansätze unterschieden, die sich hinsichtlich ihrer Natürlichkeit, des Grads der körperlichen Beteiligung und des Platzbedarfs unterscheiden \parencite{martinez2022researchTrendsInVirtualRealityLocomotionTechniques}.

\textcite{martinez2022researchTrendsInVirtualRealityLocomotionTechniques} klassifizieren Lokomotionstechniken unter anderem in selection-based, walking-based, steering-based, manipulation-based und automated Verfahren. Besonders verbreitet sind selection-based Ansätze wie Point-and-Teleport, was Nutzer:innen ermöglicht, ein Ziel im Raum auszuwählen und sich unmittelbar dorthin zu versetzen, sowie joystickbasierte Steuerungen, die eine effiziente, jedoch nur gering körperlich eingebundene Navigation ermöglichen \parencite{martinez2022researchTrendsInVirtualRealityLocomotionTechniques}. Demgegenüber zielen körperbasierte Verfahren darauf ab, die Fortbewegung enger an reale Bewegungen anzulehnen. Hierzu zählen insbesondere steering-based und walking-based Techniken, bei denen Nutzer:innen ihre Bewegungsrichtung aktiv bestimmen oder Gehbewegungen zur Navigation einsetzen.

Ein verbreiteter Ansatz für räumlich begrenzte Umgebungen ist das sogenannte Walking-in-Place, bei dem eine Gehbewegung auf der Stelle genutzt wird, um kontinuierliche Fortbewegung zu simulieren \parencite{martinez2022researchTrendsInVirtualRealityLocomotionTechniques, deBack2024naturalAndRedirectedWalkingInVirtualReality}. Darüber hinaus ermöglicht Redirected Walking eine Anpassung des Bewegungspfads durch gezielte, für Nutzer:innen kaum wahrnehmbare Verzerrungen, sodass auch in physisch begrenzten Räumen ein großräumiges Bewegungserlebnis erzeugt werden kann \parencite{deBack2024naturalAndRedirectedWalkingInVirtualReality}.

Studien zeigen, dass gehbasierte Ansätze im Vergleich zu abstrakten Methoden wie Teleportation ein stärkeres Gefühl physischer Präsenz und eine natürlichere Raumwahrnehmung fördern \parencite{martinez2022researchTrendsInVirtualRealityLocomotionTechniques, deBack2024naturalAndRedirectedWalkingInVirtualReality}. Insbesondere natürliches Gehen verbessert nachweislich die räumliche Orientierung und Distanzwahrnehmung \parencite{deBack2024naturalAndRedirectedWalkingInVirtualReality}. Gleichzeitig müssen solche Verfahren hinsichtlich Ermüdung, Tracking-Genauigkeit und physischer Sicherheit sorgfältig gestaltet werden, insbesondere in offenen Interaktionsräumen.

Für projektionsbasierte immersive Umgebungen ist Navigation daher nicht nur als technische Fortbewegung, sondern als körperlich und räumlich eingebettete Interaktionsform zu verstehen. Die Wahl eines geeigneten Lokomotionsansatzes hängt wesentlich vom verfügbaren Raum, dem Anwendungskontext und dem gewünschten Grad körperlicher Beteiligung ab.


\subsection{Gestendesign und Auswahlmechanismen}
\label{sec:rw_gestures}

Die Gestaltung geeigneter Gesten stellt einen zentralen Bestandteil controllerfreier Interaktionssysteme dar. Gesten fungieren dabei als primäres Eingabemedium und sollten sowohl intuitiv verständlich als auch technisch zuverlässig interpretierbar sein.
In der Forschung werden Gesten anhand unterschiedlicher Taxonomien klassifiziert. Eine häufig verwendete Einordnung orientiert sich an der Beziehung zwischen Geste und Bedeutung, und unterscheidet unter anderem metaphorische, symbolische, deiktische und abstrakte Gesten  \parencite{dimSilpasuwanchai2016DesigningMidAirTvGEsturesForBlindPeople}.

Metaphorische Gesten ahmen reale Handlungen nach, etwa das Greifen oder Verschieben eines Objekts, und werden häufig als besonders intuitiv wahrgenommen \parencite{kostic2024exploringmidairhandinteractionindatavisualization}.
Symbolische Gesten hingegen tragen eine konventionalisierte Bedeutung, etwa in Form bekannter Handzeichen, ohne eine reale Manipulationshandlung direkt nachzuahmen. Deiktische Gesten beziehen sich auf räumliche Informationen von Objekten, beispielsweise Zeigegesten. Abstrakte Gesten hingegen weisen keine unmittelbar nachvollziehbare Beziehung zu ihrer Funktion auf und müssen daher weitgehend erlernt werden \parencite{dimSilpasuwanchai2016DesigningMidAirTvGEsturesForBlindPeople}. Sie sind also willkürlich \parencite{dimSilpasuwanchai2016DesigningMidAirTvGEsturesForBlindPeople}.\\


Neben frei definierten Gesten greifen räumliche Interaktionssysteme häufig auf etablierte Interaktionsmuster aus anderen Domänen zurück und übertragen diese in neue Kontexte. Solche Ansätze knüpfen an bekannte Bedienlogiken an, etwa wenn Handbewegungen als Cursorsteuerung interpretiert oder Auswahlbewegungen an vertraute Zeige- und Greifhandlungen angelehnt werden. \textcite{williamsOrtega2020evolutionaryGesturesWhenAGestureIsNotQuireLegacyBiased} beschreiben solche bewussten Übertragungen als ``Evolutionary Gestures'', also Interaktionsformen, die bekannte Konzepte aufgreifen und an neue Interaktionsräume anpassen, ohne sie unverändert zu übernehmen. Die Nutzung solcher vertrauten Zuordnungen kann die Verständlichkeit und Erlernbarkeit räumlicher Interaktion unterstützen.

Gleichzeitig zeigt sich in Elicitation-Studien, in denen Nutzer:innen selbst Gestensets für bestimmte Funktionen entwickeln, dass bekannte Interaktionsmuster nicht nur gezielt durch Designer:innen übernommen, sondern auch von Nutzer:innen selbst in den dreidimensionalen Raum übertragen werden. Diese Tendenz wird als „Legacy Bias“ bezeichnet. Sie kann einerseits die Memorisierbarkeit und anfängliche Verständlichkeit von Gesten fördern, andererseits jedoch zu Interaktionsformen führen, die für räumliche Kontexte ineffizient, unergonomisch oder nur eingeschränkt geeignet sind \parencite{pasquale2024empiricalevaluationdefiningmidair}.\\


Aktuelle Forschung hebt hervor, dass es keine Standardisierung für Best Practices im Gestendesign gibt, sondern die Wahl der Geste immer kontextabhängig ist \parencite{spittle2023reviewinteractiontechniquesforimmersiveenvironments, gavgiotaki2023gesturebasedinteractionforarsystemsreview}.
Die Auswahl geeigneter Gesten erfordert deshalb die Berücksichtigung verschiedener Qualitätskriterien. Zentrale Aspekte sind dabei die Erlernbarkeit \parencite{kostic2024exploringmidairhandinteractionindatavisualization}, die sogenannte Guessability \parencite{wobbrock2005discoverability}, Intuitivität \parencite{gavgiotaki2023gesturebasedinteractionforarsystemsreview}, sowie die körperliche Belastung und soziale Akzeptanz der Gesten \parencite{mueller2010requirementsAndDesignSpaceForInteractivePublicDisplays, hincapi2014metrictoquantifyarmfatigue}. Gleichzeitig sollen Gesten so gestaltet sein, dass sie zuverlässig erkannt werden können und nicht mit alltäglichen Bewegungen verwechselt und unabsichtlich als Interaktion vom System registriert werden \parencite{kostic2024exploringmidairhandinteractionindatavisualization}.

Ein wesentliches Spannungsfeld im Gestendesign ergibt sich aus dem Zielkonflikt zwischen Natürlichkeit und Robustheit. Während stark realitätsnahe Gesten eine hohe Intuitivität aufweisen können, sind sie nicht zwangsläufig die effizientesten oder präzisesten Interaktionsformen \parencite{jacob2008realitybasedinteraction}. \\

{\color{red} TODO: Das hier eher in Menü Abschnitt? ODER Abschnitt Menübasierte Interaktion vor diesen?}
Neben der Gestaltung der Gesten selbst spielt die Auswahl- und Bestätigungsmechanik eine wichtige Rolle. Eine zentrale Herausforderung bleibt die Selektion ohne physische Bestätigung. Ohne Tasten oder Klickmechanismen fehlt haptisches Feedback, wodurch Probleme wie Overshoot, unbeabsichtigte Aktivierungen oder Unsicherheit über erfolgreiche Eingaben entstehen können.

In der Literatur werden hierfür unterschiedliche Ansätze beschrieben. Zeitbasierte Verfahren wie Dwell-Selektion aktivieren eine Auswahl automatisch, wenn ein Ziel für eine bestimmte Dauer fokussiert wird. Diese Methode erfordert keine explizite Bestätigungsgeste und wird von Nutzer:innen häufig ohne Anleitung entdeckt \parencite{walterBaillyValkanova2014CuenesicsUsingMidAirGesturesToSelectItemsOnInteractivePublicDisplays}. Allerdings kann sie zu längeren Interaktionszeiten führen \parencite{dube2022pushtabdwellpinchevaluationofselectionmethods}.

Demgegenüber stehen explizite Auswahlgesten wie Greif- oder Air-Touch-Bewegungen, die eine bewusste Bestätigung durch die Nutzer:innen erfordern. Solche Gesten orientieren sich häufig an realweltlichen Handlungen oder bekannten Interaktionsmetaphern und werden daher als besonders verständlich wahrgenommen \parencite{kostic2024exploringmidairhandinteractionindatavisualization}. Sie ermöglichen eine höhere Kontrolle über den Interaktionszeitpunkt, setzen jedoch voraus, dass die entsprechende Geste bekannt ist, korrekt ausgeführt wird und die Trackingqualität eine robuste Erkennung erlaubt \parencite{dube2022pushtabdwellpinchevaluationofselectionmethods}.
Beide Ansätze weisen unterschiedliche Vor- und Nachteile hinsichtlich Effizienz, Erlernbarkeit und Fehlerrate auf.

Insgesamt zeigt sich, dass Gestendesign und Auswahlmechanismen eng miteinander verknüpft sind und gemeinsam die Usability eines Systems bestimmen. Eine erfolgreiche Gestaltung erfordert eine sorgfältige Balance zwischen Intuitivität, Effizienz, Robustheit und ergonomischer Belastung. Diese Aspekte sind insbesondere in räumlichen Interaktionsumgebungen relevant, in denen Gesten klar unterscheidbar sein müssen und gleichzeitig nahtlos in das natürliche Verhalten der Nutzer:innen integriert werden sollen.


%Neben diskreten Auswahlmöglichkeiten spielen auch kontinuierliche Interaktionsformen eine wichtige Rolle in räumlichen Interfaces. Für die Anpassung von Parametern wie Lautstärke oder Helligkeit werden häufig Slider oder ähnliche Steuerungselemente eingesetzt. Aktuelle Ansätze kombinieren dabei Auswahlgesten mit gerichteten Bewegungen, um eine präzise und ergonomische Steuerung zu ermöglichen. Beispielsweise zeigen Systeme wie \textit{MarkAR}, dass sich Mikrogesten mit räumlichen Bewegungen kombinieren lassen, um kontinuierliche Anpassungen effizient umzusetzen \parencite{baillyPometti2024MarkARExploringTheBenefitsOfCombiningMicrogesturesAndMidAirMarksToTriggerCommandsInAR}.

\subsection{Menübasierte Interaktion}
\label{sec:rw_menus}

Zur Strukturierung von Interaktionen und zur Bereitstellung sichtbarer Funktionen spielen menübasierte Ansätze eine zentrale Rolle in räumlichen Interaktionssystemen. Insbesondere in Mid-Air- und gestenbasierten Benutzungsoberflächen adressieren Menüs ein grundlegendes Problem: die geringe Discoverability freier Gestensysteme. Da mögliche Interaktionen, die mit festgelegten Gesten erreichbar wären, häufig nicht sichtbar sind, können Nutzer:innen ohne geeignete visuelle Unterstützung Schwierigkeiten haben, verfügbare Funktionen zu erkennen und korrekt auszuführen \parencite{vogiatzidakis2018gestureeliticationstudiesreview, wobbrock2005discoverability}. Menüs lösen dieses Problem, da sie strukturiert verfügbare Funktionen visuell darstellen. \\

Räumliche Menüs übertragen etablierte Konzepte grafischer Benutzungsoberflächen in den dreidimensionalen Raum und passen diese an die Eigenschaften körperbasierter Interaktion an. Dabei kommen in der Forschung unterschiedliche Menüformen zum Einsatz, darunter beispielsweise lineare Menüs und radiale Menüs. Während lineare Menüs aus Desktop-Kontexten bekannt sind, zeigen Studien, dass diese Formen in Mid-Air-Interaktionen häufig mit erhöhtem Bewegungsaufwand und geringerer Effizienz verbunden sind, da sie längere Zielwege und präzisere Bewegungen erfordern \parencite{callahan1988comparisonPieVsLinearMenus}.
Radiale Menüs, auch als ``Pie-Menus'' bezeichnet, stellen eine besonders geeignete Alternative für räumliche Interaktion dar. Sie ordnen Funktionen kreisförmig um einen zentralen Punkt an, wodurch alle Menüeinträge in gleicher Distanz erreichbar sind. Aktuelle Vergleichsstudien in VR-Kontexten zeigen, dass radiale Menüs hohe Auswahlgeschwindigkeiten mit guter Memorierbarkeit verbinden und damit sowohl für Einsteiger:innen als auch für erfahrene Nutzer:innen geeignet sein können \parencite{wentzel2025camparisonofmenuarchetypesraycatingdirectinputmarkingmenus}.

Ein Marking‑Menu erweitert das radiale Prinzip um gestenbasierte Auswahl: Die Nutzer:innen starten im Zentrum und führen eine Richtungsbewegung aus, um das gewünschte Element zu ``markieren''. Der Auswahlpfad kann dabei mit zunehmender Übung als einzelne gerichtete Bewegung ausgeführt werden, ohne dass die visuelle Menüstruktur vollständig benötigt wird. Dies führt in Mid-Air zu kurzen, richtungsbasierten Auswahlbewegungen und reduziert sowohl den Bewegungsaufwand als auch die Fehleranfälligkeit \parencite{callahan1988comparisonPieVsLinearMenus, Mao2024comparisonOfPanelAndPieMenusForMidAirInteractionInPhotoBooth}.
Sie bieten zudem einen wichtigen Usability-Vorteil: Während Anfänger:innen auf die sichtbare Menüstruktur zurückgreifen, können erfahrene Nutzer:innen häufig verwendete Funktionen durch eingeübte Bewegungsrichtungen auswählen, ohne das Menü aktiv betrachten zu müssen. Dieser Übergang von visuell geführter zu motorisch automatisierter Interaktion wird durch das Muskelgedächtnis unterstützt und ermöglicht eine effiziente Nutzung bei wiederholten Aufgaben \parencite{kurtenbachPaul1993theDesignAndEvaluationOfMarkingMenus}. \\

Die Gestaltung von Marking-Menus wird maßgeblich durch strukturelle Parameter beeinflusst. Dazu zählen insbesondere die Anzahl der Menüeinträge pro Ebene sowie die Tiefe der Menühierarchie.
Studien deuten darauf hin, dass sowohl Genauigkeit als auch Reaktionszeit stark von der Anzahl der Einträge pro Ebene abhängen. Für Mid-Air-Pie-Menüs werden bis zu acht Einträge pro Ebene als praktikable Obergrenze beschrieben, da kleinere Zielsegmente die Auswahl erschweren \parencite{li2021depthandbreadthofpiemenusformidairgesdtureinteraction}. Gleichzeitig beeinflusst die Tiefe nur die Reaktionszeit, wobei empfohlen wird, nicht mehr als zwei Ebenen einzuführen, wenn man die optimale Interaktionszeit erzielen will \parencite{li2021depthandbreadthofpiemenusformidairgesdtureinteraction}. \\

Bestehende Menüsysteme adressieren überwiegend diskrete Auswahlaufgaben. Marking-Menus können jedoch über die rein diskrete Auswahl hinaus erweitert werden \parencite{baillyPometti2024MarkARExploringTheBenefitsOfCombiningMicrogesturesAndMidAirMarksToTriggerCommandsInAR}. Das System ``MarkAR'' zeigt beispielsweise, wie das Prinzip der gerichteten Markierung zur Steuerung kontinuierlicher und diskreter Befehle genutzt werden kann, ohne dass ein visuelles Menü explizit dargestellt werden muss. Mithilfe einer sogenannten Micro-Geste (zum Beispiel Dreifinger-Pinch) wechseln Nutzer:innen  in einen Micro-Gesten-Modus. Anstatt ein visuelles Menü anzuzeigen, interpretiert das System die Richtungsbewegung der Geste (sogenannter ``Mark'') relativ zu ihrem Startpunkt. Diese Bewegung wird auf semantische Achsen abgebildet, die als virtuelle Schieberegler zur Anpassung kontinuierlicher Werte dienen. Jedoch sind die Interaktionen in der Forschung von \textcite{baillyPometti2024MarkARExploringTheBenefitsOfCombiningMicrogesturesAndMidAirMarksToTriggerCommandsInAR} primär auf die Manipulation konkreter, funktionaler Inhalte wie Dokumente, Bilder oder 3D-Modelle ausgelegt und weniger auf die Einstellung beliebiger, kontinuierlicher Parameter ausgerichtet. Weiterhin war eine Limitation der Kompromiss zwischen Genauigkeit und Gestenamplitude, denn Teilnehmer der Studie berichteten, dass es schwierig war, präzise Werte anzusteuern (zum Beispiel eine bestimmte Seite in einem PDF). Damit priorisierte das System kleine, ermüdungsarme Gesten, was jedoch die motorische Genauigkeit einschränkte.

Insgesamt verdeutlicht die Literatur, dass räumliche Menüs eine zentrale Rolle bei der Gestaltung nutzbarer Mid-Air-Interaktionssysteme spielen. Sie ermöglichen die Sichtbarmachung von Funktionen, strukturieren komplexe Interaktionsräume und unterstützen sowohl Einsteiger:innen als auch erfahrene Nutzer:innen. Insbesondere radiale Menüs bieten durch ihre gleichmäßige Anordnung, ihre Eignung für richtungsbasierte Interaktion sowie ihre Unterstützung von motorischem Lernen entscheidende Vorteile für den Einsatz in räumlichen und körperbasierten Interaktionskontexten. Gleichzeitig zeigt sich, dass die Integration kontinuierlicher Steuerungsformen in solche menübasierten Konzepte bisher nur begrenzt untersucht wurde.


\subsection{Multisensorische und atmosphärische Interaktion}
\label{sec:rw_multisensory}
Aktuelle Interaktionssysteme adressieren häufig nur einen begrenzten Teil der menschlichen Sinneswahrnehmung, wobei visuelle und auditive Reize dominieren, während andere Modalitäten wie Berührung oder Geruch oft unberücksichtigt bleiben. Multisensorische immersive Umgebungen verfolgen das Ziel, diese Einschränkung zu überwinden, indem sie mehrere Sinneskanäle gleichzeitig ansprechen und so ein umfassenderes Erleben digitaler Inhalte ermöglichen.

Solche Systeme werden in unterschiedlichen Kontexten erforscht und eingesetzt, etwa in Bereichen wie Therapie, Unterhaltung oder erfahrungsorientierten Ausstellungs- und Interaktionsräumen. Ein Beispiel hierfür ist das von \textcite{matar2023designOfAMixedRealityBasedImmersiveVirtualEnvironmentSystem} vorgestellte System, das eine 270°-Umgebung mit wandfüllender Projektion nutzt und Interaktion über Körperposition und Handgesten ermöglicht. Nutzer:innen können dabei durch ihre Bewegung im Raum sowie durch Handgesten virtuelle Umgebungen manipulieren. Im Unterschied zu individuell getragenen XR-Systemen steht dabei eine gemeinsam nutzbare, nicht wearable-basierte räumliche Interaktion im Vordergrund. Das Beispiel verdeutlicht, dass projektionsbasierte immersive Umgebungen nicht auf rein rezeptive Erfahrungen beschränkt sind, sondern auch aktive Eingriffe in die dargestellte Szene unterstützen können. Gleichzeitig wird deutlich, dass solche Systeme mit technischen und konzeptionellen Herausforderungen verbunden sind, insbesondere hinsichtlich der präzisen Gestenerkennung, der Synchronisation von Audio- und Videokomponenten sowie der Integration in bestehende räumliche Infrastrukturen.

Der Multisensorische Interaktionsraum (MIR) lässt sich in diesen weiteren Kontext räumlicher, projektionsbasierter und multisensorischer Systeme einordnen. Er stellt eine Infrastruktur bereit, die 270°-Projektion, Mehrkanal-Audio, Beduftung und taktile Displays kombiniert, um immersive Erlebnisse zu erzeugen \parencite{orawetz2024mir}. Im Unterschied zu klassischen Interaktionssystemen steht dabei nicht ausschließlich die effiziente Ausführung einzelner Aufgaben im Vordergrund. Je nach Szenario kann der Fokus vielmehr auf der Gestaltung von Atmosphäre, dem Erleben eines narrativen oder räumlichen Settings oder auch auf gezielten funktionalen Eingriffen in das System liegen.

Daraus ergeben sich besondere Anforderungen an die Interaktionsgestaltung. Interaktion in multisensorischen Umgebungen muss nicht nur funktional sein, sondern sich in Wahrnehmung, Raum und Nutzungskontext einfügen. Je nach Anwendung kann es sinnvoll sein, Interaktion möglichst unaufdringlich zu gestalten, um das Erleben einer Szene nicht unnötig zu unterbrechen. In anderen Kontexten sind dagegen explizite und deutlich sichtbare Eingriffe erforderlich, etwa wenn Inhalte gezielt gesteuert oder verändert werden sollen. Entscheidend ist daher weniger die vollständige Unsichtbarkeit von Interaktion als vielmehr ihre Kontextangemessenheit.

Designmodelle wie das von \textcite{vieria2024understandingAffordancesinXR} vorgeschlagene Modell zur Integration multimodaler Interaktion betonen, dass die Abstimmung verschiedener Modalitäten wesentlich für glaubwürdige und interaktive immersive Umgebungen ist. Die Kombination aus visuellen, auditiven, haptischen und weiteren sensorischen Rückmeldungen kann das Erleben von Präsenz und Kontrolle unterstützen, erfordert jedoch eine kohärente Gestaltung der Interaktionsformen und ihrer Rückmeldungen \parencite{vieria2024understandingAffordancesinXR}.

Ein Beispiel für die Integration zielgerichteter Gesteninteraktion im Kontext des MIR stellt AVISAR dar, bei dem Gesten zur Steuerung von Reparaturanleitungen innerhalb einer immersiven Projektionsumgebung eingesetzt werden \parencite{stoll2024avisar, natschke2024display}. Dieses Beispiel zeigt, dass in multisensorischen und immersiven Räumen nicht nur atmosphärische, sondern auch aufgabenbezogene Interaktionen relevant sein können.


Insgesamt zeigt der Forschungsstand, dass zwar zahlreiche Ansätze für körperbasierte, gestenbasierte und räumliche Interaktion in immersiven Umgebungen vorliegen, dass jedoch insbesondere für offene, projektionsbasierte und multisensorische Räume weiterhin Forschungsbedarf besteht. Dies betrifft vor allem Interaktionskonzepte, die controllerfrei, szenariounabhängig, ergonomisch und zugleich mit atmosphärischen wie funktionalen Nutzungskontexten vereinbar sein sollen.


%Bisher existieren wenige Ansätze, die controllerfreie, körperbasierte Interaktionen in multisensorischen, projektionsbasierten Räumen wie dem MIR mit der kontinuierlichen, atmosphärischen Steuerung von Umgebungsparametern verbinden. Während Marking-Menus in VR- oder AR-Kontexten untersucht wurden, fehlen systematische Konzepte für ihre Anwendung in offenen, gemeinsam genutzten, multisensorischen Räumen, in denen Nutzer:innen nicht nur Funktionen auswählen, sondern die Atmosphäre eines Raumes dynamisch beeinflussen sollen. Insbesondere die Integration von richtungsbasierten Gesten mit kontinuierlichen, multimodalen Rückmeldungen (zum Beispiel Geruch, Temperatur, Taktik) bleibt unerforscht. Diese Arbeit adressiert diese Lücke, indem sie ein gestengesteuertes Marking Menu für den MIR konzipiert und prototypisch umsetzt, das sowohl für atmosphärische als auch funktionale Interaktionen geeignet ist.

